\documentclass[10pt,a4paper]{amsart}
\usepackage[margin=19mm]{geometry}
\usepackage{amsmath,amssymb,mathtools}
\usepackage[scr=rsfs,cal=euler]{mathalfa}
\usepackage{xfrac}
\usepackage[unicode,hidelinks,bookmarks=false]{hyperref}
\newcommand{\ie}{i.e.\ }
\newcommand{\cf}{cf.\ }
\newcommand{\resp}{respectively\ }
\newcommand{\Subsection}{\subsection}
\providecommand{\nobreakdash}{\nobreak\hbox{-}}
\setlength{\emergencystretch}{2em}
\multlinegap0pt
\usepackage{bm}
\usepackage{bbm}
\usepackage{dsfont}
\usepackage{xspace}
\usepackage{cancel}
\usepackage{mathtools}
\usepackage{amsthm}
\makeatletter
\@ifundefined{theorem}{\newtheorem{theorem}{Theorem}}{}
\@ifundefined{proposition}{\newtheorem{proposition}[theorem]{Proposition}}{}
\makeatother
\def\dif{{\mathord{{\rm d}}}}
\def\geq{\geqslant}
\def\leq{\leqslant}
\def\min{{\mathord{{\rm min}}}}
\title[The Wave Kinetic Theory for Quasilinear MMT Equation]{The Wave Kinetic Theory for Quasilinear MMT Equation}
\author{Huaxiang L\"u}
\date{Source: arXiv:2608.22710v1 (CC BY 4.0)}
\begin{document}
\maketitle

\noindent\emph{Source and licence.} The Wave Kinetic Theory for Quasilinear MMT Equation. Version arXiv:2608.22710v1 (CC BY 4.0), distributed under the Creative Commons Attribution 4.0 International licence, as recorded in the arXiv licence metadata for this exact version and shown on the abstract page https://arxiv.org/abs/2608.22710v1. Full rights notice: licenses/arxiv-2608.22710-CC-BY-4.0.txt.\par\smallskip

\noindent\textit{Editorial scope.} Counted unit: the Proposition whose source label is \texttt{asymptotic} in the article (source0.tex lines 1489--1502, source label \texttt{asymptotic}), delivered together with the article's own complete original proof of it (source0.tex lines 1504--1576, one \texttt{proof} environment of 73 lines). No mathematics is written, repaired or re-derived: statement and proof are the article's own text line for line. Required context retained uncounted: none. Every definition, display and label of those lines is retained unchanged. Omitted from this package and not redistributed: 1--1488, 1503--1503, 1577--5692. Nothing inside a retained excerpt is omitted or altered: every retained body is delivered exactly as the article's lines are written, so no omission receipt of any kind accompanies this package. Citations: the retained excerpts carry no citation command at all, so no bibliography entry of the article is transcribed and no reference is redistributed. Reference adaptation: none was needed and none was made; every reference inside the delivered unit resolves inside it, so no unresolved reference or citation is produced. Printed slips kept literal: no printed word, symbol, hypothesis or number of the counted statement or of its proof was altered. Layout and class: the article's own class is \texttt{amsart} with packages including enumerate, mhequ, amsthm, amsmath, amssymb, url, color, booktabs, nccmath, geometry, mathrsfs, enumitem, dsfont, subfigure; the wrapper is the lane's pinned \texttt{amsart} editorial wrapper at 10pt on A4, which restates the theorem-like environments the retained text needs and adds the article's own macro definitions that the retained text needs (\texttt{\textbackslash dif}, \texttt{\textbackslash geq}, \texttt{\textbackslash leq}, \texttt{\textbackslash min}), with extra wrapper packages (bm, bbm, dsfont, xspace, cancel, mathtools). The delivered numbering is the unit's own. Assets: no retained span contains a figure environment, a graphics-inclusion command, a PDF-page inclusion, a table or a TikZ picture; no image file is copied into this package, the package declares no asset dependency and no image file is redistributed with it. Licence: version arXiv:2608.22710v1 of this article is distributed under the Creative Commons Attribution 4.0 International licence, as recorded in the arXiv licence metadata for this exact version and shown on the abstract page https://arxiv.org/abs/2608.22710v1. A full rights notice is delivered at licenses/arxiv-2608.22710-CC-BY-4.0.txt. Characters: every retained excerpt is pure ASCII, so the delivered engine, XeLaTeX with its default Latin Modern text font, reports no missing character.
\begin{proposition}
    \label{asymptotic}
Let $T  < L^{\frac{(1+\gamma)(1+2\beta)}{2}-10\delta}$ if $\sigma\in(1,2]$, $T  <L^{\frac{1+2\beta}{2+2\beta-\sigma}-10\delta} $ if $\sigma\in(0,1)$, and $\phi \in \mathcal{S}(\mathbb{R})$. 
For any $t \in (0,T)$, let   \begin{align*}
 \mathscr S_t(\phi):&=\sum\limits_{\substack{ k-k_1+k_2-k_3=0,\\ \overline \Omega_k\neq0}}\phi_k \phi_{k_1} \phi_{k_2} \phi_{k_3}  \left[ \frac{1}{\phi_k} - \frac{1}{\phi_{k_1}} + \frac{1}{\phi_{k_2}} - \frac{1}{\phi_{k_3}} \right]\left| \frac{\sin(\pi t\overline \Omega_k)}{\pi t \overline \Omega_k} \right|^2\notag\\
  &\quad\quad\quad\times
|k_1|^{2\beta}|k_2|^{2\beta}|k_3|^{2\beta}|k|^{2\beta}\varphi_{\leq K}^2(k_1)\varphi_{\leq K}^2(k_2)\varphi_{\leq K}^2(k_3)\varphi_{\leq K}^2(k).\\
\mathscr K_t(\phi):&=L^{2} \int_{\xi_1-\xi_2+\xi_3=\xi }\phi(\xi) \phi(\xi_1) \phi(\xi_2) \phi(\xi_3)  \left[ \frac{1}{\phi(\xi)} - \frac{1}{\phi(\xi_1)} + \frac{1}{\phi(\xi_2)} - \frac{1}{\phi(\xi_3)} \right]    \left| \frac{\sin(\pi t\overline \Omega( \xi))}{\pi t\overline\Omega( \xi)} \right|^2 \notag\\ &\quad\quad\times |\xi_1|^{2\beta}|\xi_2|^{2\beta}|\xi_3|^{2\beta}|\xi|^{2\beta}\varphi_{\leq K}^2(\xi_1)\varphi_{\leq K}^2(\xi_2)\varphi_{\leq K}^2(\xi_3)\varphi_{\leq K}^2(\xi)   \dif \xi_1  \dif\xi_2 \, \dif\xi_3.
\end{align*}
 Then
$$
\mathscr S_t=\mathscr K_t +O(L^{2-\delta} t^{-1}).
$$
\end{proposition}

\begin{proof}
First, fix $k$ and define \begin{align*}
    g(x):&=\left|\frac{\sin \pi x}{\pi x}\right|^2\leq 1,\\
    F(k_1,k_3):&=\phi_k \phi_{k_1} \phi_{k_1+k_3-k} \phi_{k_3}  \left[ \frac{1}{\phi_k} - \frac{1}{\phi_{k_1}} + \frac{1}{\phi_{k_1+k_3-k}} - \frac{1}{\phi_{k_3}} \right]\\
    &\quad\quad\times
 |k_1|^{2\beta}|k_1+k_3-k|^{2\beta}|k_3|^{2\beta}|k|^{2\beta}\varphi_{\leq K}^2(k_1)\varphi_{\leq K}^2(k_1+k_3-k)\varphi_{\leq K}^2(k_3)\varphi_{\leq K}^2(k)\\
 &\quad\quad \times \chi(T^{\frac{1}{1+2\beta}}L^{2\delta}|k_1|)\chi(T^{\frac{1}{1+2\beta}}L^{2\delta}|k_1+k_3-k|)\chi(T^{\frac{1}{1+2\beta}}L^{2\delta}|k_3|),
\end{align*}
where   $\chi$ is a smooth cutoff satisfying $\chi(x)=0$ for $|x|\leq \frac12$ and $\chi(x)=1$ for $|x|\geq 1$. These truncations remove the singularity near $|k_i|=0$. 

Define the truncated quantities
\begin{align*}
   (\mathscr S_{\rm tr})_t&:=\sum_{(k_1, k_3)\in \mathbb{Z}_L^{2},\overline\Omega(k_1,k_3)\neq 0}F(k_1,k_3) g\left(t\overline \Omega(k_1,k_3)\right),\\
 (\mathscr K_{\rm tr})_t&:=L^{2} \int_{\xi_1,\xi_3}F(\xi_1,\xi_3) g\left(t\overline \Omega(\xi_1,\xi_3)\right)\dif \xi_1\dif \xi_3,
\end{align*}
where
$$\overline\Omega(x,y) =|x|^\sigma-|x+y-k|^\sigma+|y|^\sigma-|k|^\sigma+\Gamma_0(x)-\Gamma_0(x+y-k)+\Gamma_0(y)-\Gamma_0(k),$$
and
$ \Gamma_0(x)=C_K\frac{\alpha }{\pi }\varphi_{\leq K}(x)|x|^{2\beta}.$

In the first step,  for  the truncated asymptotic functions, we show that $$
( \mathscr S_{\rm tr})_t= (\mathscr K_{\rm tr})_t +O(L^{2-\delta} t^{-1}).
$$
Since $\widehat g(\tau)=1-|\tau|$ on $[-1,1]$ and vanishes outside, we write $e(x):=e^{2\pi i x}$. Applying Poisson summation, we obtain 
\begin{align*}
 (  \mathscr S_{\rm tr})_t&=\sum_{(k_1, k_3)}F(k_1,k_3) g\left(t\overline \Omega(k_1,k_3)\right)=\sum_{(k_1, k_3)}F(k_1,k_3)\int_{-\infty}^{\infty}\hat{g}(\tau)e(\tau t\overline \Omega(k_1,k_3))\dif \tau\\
 &=\frac1t \sum_{(k_1, k_3)}F(k_1,k_3)\int_{-\infty}^{\infty}\hat{g}(\frac\tau t)e(\tau \overline \Omega(k_1,k_3))\dif \tau=\frac1t \int_{-\infty}^{\infty}\hat{g}(\frac\tau t)\sum_{(k_1, k_3)}F(k_1,k_3)e(\tau \overline \Omega(k_1,k_3))\dif \tau\\
 &=\frac1t \int_{-\infty}^{\infty}\hat{g}(\frac\tau t)\sum_{(k_1, k_3)\in\mathbb{Z}^{2}}F(\frac{k_1}{L},\frac{k_3}L)e(\tau \overline \Omega(\frac{k_1}L,\frac{k_3}L))\dif \tau\\
 &=\frac1t \int_{-\infty}^{\infty}\hat{g}(\frac\tau t) \sum_{(c,d)\in\mathbb{Z}^{2}}\int_{(x,y)\in\mathbb{R}^{2}}F(\frac{x}{L},\frac{y}L)e(\tau \overline \Omega(\frac{x}L,\frac{y}L))e(-c\cdot x-d\cdot y) \dif x\dif y\dif \tau\\
  &=\frac1t \int_{-\infty}^{\infty}\hat{g}(\frac\tau t) \int_{(x,y)\in\mathbb{R}^{2}}F(\frac{x}{L},\frac{y}L)e(\tau \overline \Omega(\frac{x}L,\frac{y}L))\dif x\dif y\dif \tau=:  (\mathscr S_{\rm tr})_1\\
   &\quad+\frac{L^{2}}t \int_{-\infty}^{\infty}\hat{g}(\frac\tau t) \sum_{(c,d)\in\mathbb{Z}^{2}/\{0,0\}}\int_{(x,y)\in\mathbb{R}^{2}}F(x,y)e(\tau \overline \Omega(x,y)-c\cdot Lx-d\cdot Ly) \dif x\dif y\dif \tau=:  (\mathscr S_{\rm tr})_2.
\end{align*}

For the main term, a change of variables yields
\begin{align*}
  ( \mathscr S_{\rm tr})_1=\frac{L^{2}}t \int_{-\infty}^{\infty}\hat{g}(\frac\tau t) \int_{(x,y)\in\mathbb{R}^{2}}F(x,y)e(\tau \overline \Omega(x,y))\dif x\dif y\dif \tau= (\mathscr K_{\rm tr})_t.
\end{align*}

For the second term $(\mathscr S_{\rm tr})_2$, in the case $\sigma\in(1,2]$, it holds that by the cut off function $\varphi$ and $\chi$,
\begin{align*}
   |\tau  \nabla  \overline \Omega(x,y)|\lesssim T(L^{\delta(\sigma-1)}+\alpha (T^{\frac{1}{1+2\beta}}L^{2\delta})^{1-2\beta})\leq L^{1-\delta/2},
\end{align*}
where we used $|\tau|\leq T < \min\{L^{1-10\delta},L^{\frac{(1+\gamma)(1+2\beta)}{2}-10\delta}\}$.

While in the case $\sigma\in(0,1)$, we have 
\begin{align*}
   |\tau  \nabla  \overline \Omega(x,y)|\lesssim T((T^{\frac{1}{1+2\beta}}L^{2\delta})^{1-\sigma}+\alpha (T^{\frac{1}{1+2\beta}}L^{2\delta})^{1-2\beta})\leq L^{1-\delta/2},
\end{align*}
where we used $|\tau|\leq T < L^{\frac{1+2\beta}{2+2\beta-\sigma}-10\delta}$ and $2\beta<\sigma$. In summary, we have 
\begin{align*}
   |\tau  \nabla  \overline \Omega(x,y)| \leq L^{1-\delta/2}.
\end{align*}

Then for $c\neq 0$ we have
\begin{align*}
|   \nabla_x[ \tau \overline\Omega(x,y)-c\cdot Lx-d\cdot Ly]|\geq L|c|-CL^{1-\delta/2}\geq \frac12|c|L,
\end{align*}
as $L$ goes large enough.


Now, from the definition of $F(x,y)$,  once we apply integration by parts, 
we notice that by the cut off function $\varphi$ and $\chi$,  $$|\partial_{x_i}|x|^{2\beta}|\leq 2\beta|x|^{2\beta-1}\lesssim (TL^{2\delta})^{1-2\beta},$$ 
we obtain a power of $TL^{2\delta}\cdot (|c|L)^{-1}$ from $F(x,y)$. By  integrating enough times, we have a gain of $L^{-2\delta}$,  then we obtain $(\mathscr S_{\rm tr})_2 =O(L^{2-\delta} t^{-1}).$

It remains to estimate the difference between truncated and original quantities. By the support of $1-\chi$, at least one of $|k_1|,|k_1+k_3-k|,|k_3|$ is $\lesssim T^{-{\frac{1}{1+2\beta}}}L^{-2\delta}$, and hence
\begin{align*}
 | ( \mathscr K_{\rm tr})_t- \mathscr K_t|&\leq   L^{2} \int_{\xi_1,\xi_3}\phi_\xi \phi_{\xi_1} \phi_{\xi_1+\xi_3-\xi} \phi_{\xi_3}  \left[ \frac{1}{\phi_\xi} - \frac{1}{\phi_{\xi_1}} + \frac{1}{\phi_{\xi_1+\xi_3-\xi}} - \frac{1}{\phi_{\xi_3}} \right](T^{-{\frac{1}{1+2\beta}}}L^{-2\delta} L^{\delta})^{1+2\beta}\dif \xi_1\dif \xi_3\\
 &=O(L^2(T^{-{\frac{1}{1+2\beta}}}L^{-\delta})^{1+2\beta})\leq O(L^{2-\delta}T^{-1})\leq O(L^{2-\delta}t^{-1}).
\end{align*} The same argument yields
\begin{align*}
 | (\mathscr S_{\rm tr})_t- \mathscr S_t|&\leq  O(L^{2-\delta}t^{-1}).
\end{align*}Combining the above estimates, we conclude the proof.
\end{proof}

\end{document}
